% ============================================================================
% LANGENTWURF V2: Mi instituto - Objetos de clase & hay
% Klasse 7 | Spanisch | Doppelstunde (90 Min.)
% ============================================================================
% VERSION 2 - Neuer Workflow
% Lehrwerk: Qué pasa 1, Unidad 2 (S. 36ff.)
% Fokus: Klassenzimmer-Objekte + hay
% NICHT: estar, Lektionstext, -er/-ir Verben
% ============================================================================

\documentclass[11pt,a4paper]{article}

% ============================================================================
% PACKAGES
% ============================================================================
\usepackage[utf8]{inputenc}
\usepackage[T1]{fontenc}
\usepackage[ngerman]{babel}
\usepackage{geometry}
\usepackage{fancyhdr}
\usepackage{lastpage}
\usepackage{xcolor}
\usepackage{tcolorbox}
\usepackage{enumitem}
\usepackage{tabularx}
\usepackage{longtable}
\usepackage{array}
\usepackage{booktabs}
\usepackage{hyperref}
\usepackage{amssymb}

% ============================================================================
% KONFIGURATION
% ============================================================================

\newcommand{\plantier}{1}
\newcommand{\fachid}{spanisch}
\newcommand{\fach}{Spanisch}
\newcommand{\fachfarbe}{c41e3a}

% ============================================================================
% VARIABLEN
% ============================================================================

\newcommand{\jahrgangsstufe}{7}
\newcommand{\schulart}{Gesamtschule}
\newcommand{\stundenthema}{Objetos de clase -- Was gibt es im Klassenraum?}
\newcommand{\reihenthema}{Unidad 2: Mi instituto}
\newcommand{\stundennummer}{1}
\newcommand{\reihenstunden}{4}
\newcommand{\unterrichtsdatum}{\_\_\_\_\_\_\_\_\_\_}
\newcommand{\unterrichtszeit}{\_\_\_\_\_ -- \_\_\_\_\_}
\newcommand{\raum}{\_\_\_\_\_}

\newcommand{\lehrkraft}{N.N.}
\newcommand{\ausbildungsstand}{Lehrkraft}

\newcommand{\lerngruppengroesse}{24}
\newcommand{\maennlich}{12}
\newcommand{\weiblich}{12}
\newcommand{\divers}{0}

\newcommand{\gerniveau}{A1}
\newcommand{\lehrwerk}{Qué pasa 1}
\newcommand{\lehrwerkseite}{S. 36--38}

% ============================================================================
% LAYOUT
% ============================================================================
\geometry{left=2.5cm, right=2.5cm, top=2.5cm, bottom=2.5cm}
\definecolor{fach}{HTML}{\fachfarbe}
\definecolor{gray}{HTML}{666666}
\definecolor{lightgray}{HTML}{f5f5f5}
\definecolor{successgreen}{HTML}{2a7a4b}
\definecolor{warningorange}{HTML}{b35c00}

\tcbuselibrary{skins,breakable}

\newtcolorbox{infobox}[1][]{
    colback=lightgray,
    colframe=fach,
    fonttitle=\bfseries,
    title=#1,
    breakable
}

\newtcolorbox{scorebox}{
    colback=white,
    colframe=fach,
    fonttitle=\bfseries,
    title=Didaktik-Score,
    breakable
}

\pagestyle{fancy}
\fancyhf{}
\fancyhead[L]{\small\textcolor{gray}{\fach{} -- Jg. \jahrgangsstufe{} -- \stundenthema}}
\fancyhead[R]{\small\textcolor{gray}{Langentwurf Tier 1 -- V2}}
\fancyfoot[C]{\small\thepage{} / \pageref{LastPage}}
\renewcommand{\headrulewidth}{0.4pt}
\renewcommand{\footrulewidth}{0pt}

% ============================================================================
% DOKUMENT
% ============================================================================
\begin{document}

% ============================================================================
% DECKBLATT
% ============================================================================
\begin{titlepage}
    \centering
    \vspace*{2cm}

    {\large\textcolor{gray}{\schulart{} -- Version 2}}\\[2cm]

    {\Huge\textcolor{fach}{\textbf{Langentwurf}}}\\[0.5cm]
    {\large Tier 1 -- Vollständig}\\[2cm]

    {\LARGE\textbf{\stundenthema}}\\[0.5cm]
    {\large im Rahmen der Unterrichtsreihe}\\[0.3cm]
    {\Large\textit{\reihenthema}}\\[2cm]

    \begin{tabular}{rl}
        \textbf{Fach:} & \fach{} (\gerniveau) \\[0.3cm]
        \textbf{Klasse:} & \jahrgangsstufe \\[0.3cm]
        \textbf{Lehrwerk:} & \lehrwerk, \lehrwerkseite \\[0.3cm]
        \textbf{Datum:} & \underline{\hspace{3cm}} \\[0.3cm]
        \textbf{Zeit:} & \underline{\hspace{3cm}} (Doppelstunde) \\[0.3cm]
        \textbf{Raum:} & \underline{\hspace{3cm}} \\[0.3cm]
    \end{tabular}

    \vfill

    {\large Lehrkraft: \lehrkraft}\\[0.3cm]
    {\normalsize\textcolor{gray}{\ausbildungsstand}}\\[1cm]

\end{titlepage}

% ============================================================================
% INHALTSVERZEICHNIS
% ============================================================================
\tableofcontents
\newpage

% ============================================================================
% 1. BEDINGUNGSANALYSE
% ============================================================================
\section{Bedingungsanalyse}

\subsection{Lerngruppe}
\begin{tabular}{ll}
    Kursgröße: & \lerngruppengroesse{} SuS (\maennlich{} m / \weiblich{} w / \divers{} d) \\
    Jahrgangsstufe: & \jahrgangsstufe \\
    Sprachniveau: & \gerniveau{} (GER) \\
    Fremdsprache: & 2. Fremdsprache (1. Lernjahr) \\
\end{tabular}

\subsection{Vorkenntnisse (Unidad 0--1)}
\begin{itemize}
    \item Begrüßung, Verabschiedung, Befinden
    \item Verb \textit{ser} (yo soy, tú eres, él/ella es...)
    \item Verb \textit{tener} (tengo, tienes, tiene...)
    \item Bestimmter Artikel: el, la, los, las
    \item Unbestimmter Artikel: un, una
    \item Zahlen 0--20
    \item Erste -ar Verben (hablar, estudiar)
\end{itemize}

\subsection{Differenzierung}

\textbf{Hinweis:} Keine Niveaumarkierung auf Schülermaterial!

\begin{tabular}{|l|p{3.8cm}|p{3.8cm}|p{3.8cm}|}
\hline
& \textbf{[B] Basis} & \textbf{[S] Standard} & \textbf{[E] Erweiterung} \\
\hline
\textbf{Zielgruppe} & Förderbedarf, LRS & Regelschüler & Leistungsstark \\
\hline
\textbf{Vokabeln} & 8 Kernwörter (rezeptiv) & 12--14 Wörter (produktiv) & 18+ Wörter (Transfer) \\
\hline
\textbf{Grammatik} & \textit{hay} + Bild erkennen & \textit{hay} + Sätze bilden & \textit{hay} + freie Beschreibung \\
\hline
\textbf{Output} & Zeigen, Nachsprechen & Lückensätze, Kurzantworten & Eigene Sätze, Mini-Text \\
\hline
\end{tabular}

\subsection{Besonderheiten der Lerngruppe}
\begin{itemize}
    \item \textbf{LRS:} 2 SuS $\rightarrow$ Vergrößerte Schrift, Zeitzugabe
    \item \textbf{DaZ:} 1 SuS $\rightarrow$ Deutsch-Spanisch parallel einführen
    \item \textbf{Hochbegabt:} 1 SuS $\rightarrow$ Expertenrolle, Zusatzaufgaben
\end{itemize}

% ============================================================================
% 2. SACHANALYSE
% ============================================================================
\section{Sachanalyse}

\subsection{Wortfeld: Objetos de clase}

\subsubsection{Kernvokabular (für alle)}

\begin{center}
\begin{tabular}{|l|l|l|}
\hline
\textbf{Spanisch} & \textbf{Deutsch} & \textbf{Artikel-Tipp} \\
\hline
la mesa & der Tisch & fem. \\
la silla & der Stuhl & fem. \\
la pizarra & die Tafel & fem. \\
la ventana & das Fenster & fem. \\
la puerta & die Tür & fem. \\
el libro & das Buch & mask. \\
el cuaderno & das Heft & mask. \\
el bolígrafo (boli) & der Kugelschreiber & mask. \\
\hline
\end{tabular}
\end{center}

\subsubsection{Erweiterungsvokabular}

\begin{center}
\begin{tabular}{|l|l|l|}
\hline
\textbf{Spanisch} & \textbf{Deutsch} & \textbf{Hinweis} \\
\hline
la mochila & der Rucksack & -- \\
el lápiz & der Bleistift & Plural: los lápices \\
el estuche & das Mäppchen & -- \\
la carpeta & die Mappe & -- \\
la papelera & der Papierkorb & -- \\
el mapa & die Landkarte & \textbf{Achtung: maskulin!} \\
la pared & die Wand & -- \\
el reloj & die Uhr & -- \\
la calculadora & der Taschenrechner & -- \\
el ordenador & der Computer & -- \\
\hline
\end{tabular}
\end{center}

\textbf{Besonderheit:} \textit{el mapa} -- maskulin trotz Endung -a (griechischer Ursprung).

\subsection{Grammatik: hay (es gibt)}

\subsubsection{Form und Bedeutung}

\textbf{hay} ist die unpersönliche Form von \textit{haber} im Präsens.
\begin{itemize}
    \item Bedeutung: ``es gibt'' (Singular und Plural!)
    \item Unveränderlich: hay bleibt immer hay
\end{itemize}

\subsubsection{Verwendung}

\begin{center}
\begin{tabular}{|l|l|}
\hline
\textbf{Struktur} & \textbf{Beispiel} \\
\hline
hay + un/una + Nomen & Hay \textbf{una} mesa. \\
hay + Zahl + Nomen & Hay \textbf{cinco} sillas. \\
hay + muchos/as + Nomen & Hay \textbf{muchos} libros. \\
hay + Nomen (Plural ohne Artikel) & Hay cuadernos. \\
\hline
\end{tabular}
\end{center}

\subsubsection{Typische Fehler}

\begin{center}
\begin{tabular}{|l|l|l|}
\hline
\textbf{Fehler} & \textbf{Korrektur} & \textbf{Erklärung} \\
\hline
*Hay el libro & Hay \textbf{un} libro & Kein bestimmter Artikel nach hay! \\
*Hay la mesa & Hay \textbf{una} mesa & hay = unbestimmt/neu \\
*Es hay una mesa & \textbf{Hay} una mesa & Kein Subjekt vor hay \\
\hline
\end{tabular}
\end{center}

\subsubsection{Fragen mit hay}

\begin{itemize}
    \item \textit{¿Qué hay en la clase?} -- Was gibt es in der Klasse?
    \item \textit{¿Cuántos libros hay?} -- Wie viele Bücher gibt es?
    \item \textit{¿Cuántas sillas hay?} -- Wie viele Stühle gibt es?
\end{itemize}

\subsection{Kontrastive Analyse}

\begin{tabular}{|p{6cm}|p{7cm}|}
\hline
\textbf{Transfer (positiv)} & \textbf{Interferenz (negativ)} \\
\hline
``Es gibt'' ist auch im Deutschen unpersönlich & Tendenz zum bestimmten Artikel (*hay el...) \\
\hline
Wortstellung ähnlich: hay + Objekt & Genus-Unterschiede (el mapa, la mano) \\
\hline
Unveränderlichkeit leicht zu merken & Verwechslung mit \textit{estar} (kommt später) \\
\hline
\end{tabular}

% ============================================================================
% 3. DIDAKTISCHE ANALYSE
% ============================================================================
\section{Didaktische Analyse}

\subsection{Einbettung in die Sequenz}

\textbf{Sequenz ``Mi instituto'' (4 Doppelstunden):}

\begin{enumerate}
    \item \textbf{DS 1 (diese Stunde):} Objetos de clase + hay
    \item DS 2: estar + Ortsangaben (¿Dónde está...?)
    \item DS 3: Verben auf -er/-ir + Schulalltag
    \item DS 4: Repaso + Leistungsüberprüfung
\end{enumerate}

\subsection{Kompetenzen (Rahmenplan MV)}

\begin{tabular}{|l|p{10cm}|}
\hline
\textbf{Bereich} & \textbf{Kompetenz} \\
\hline
Sprechen & Einfache Aussagen über den unmittelbaren Lebensbereich machen \\
\hline
Wortschatz & Grundwortschatz Schule/Klassenraum aufbauen \\
\hline
Grammatik & Unpersönliche Konstruktion \textit{hay} verstehen und anwenden \\
\hline
Interkulturell & Spanischen Schulalltag kennenlernen (instituto) \\
\hline
\end{tabular}

\subsection{Didaktische Reduktion}

\textbf{Bewusst ausgeklammert in dieser Stunde:}
\begin{itemize}
    \item \textit{estar} (Ortsangaben) $\rightarrow$ DS 2
    \item Lektionstext (S. 38--39) $\rightarrow$ DS 2
    \item -er/-ir Verben $\rightarrow$ DS 3
    \item Negation \textit{no hay} $\rightarrow$ implizit, nicht fokussiert
\end{itemize}

\textbf{Begründung:} Fokussierung auf einen grammatischen Schwerpunkt ermöglicht tieferes Verständnis und mehr Übungszeit.

% ============================================================================
% 4. STUNDENZIELE
% ============================================================================
\section{Stundenziele}

\subsection{Grobziel}

Die SuS erweitern ihre \textbf{kommunikative Kompetenz im Bereich Sprechen}, indem sie Gegenstände im Klassenraum benennen und mit \textit{hay} beschreiben, was es gibt.

\subsection{Feinziele}

\begin{tabular}{|c|p{7.5cm}|c|c|c|}
\hline
\textbf{FZ} & \textbf{Die SuS können...} & \textbf{AFB} & \textbf{Diff.} & \textbf{Phase} \\
\hline
1 & 8 Klassenzimmer-Vokabeln korrekt benennen & I & alle & St. 1 \\
\hline
2 & den Artikel (el/la) zu Vokabeln zuordnen & I & alle & St. 1 \\
\hline
3 & \textit{hay} + un/una in Aussagesätzen verwenden & II & alle & St. 1 \\
\hline
4 & die Frage \textit{¿Qué hay en...?} verstehen und beantworten & II & [S]/[E] & St. 2 \\
\hline
5 & \textit{hay} + Zahl zur Mengenangabe nutzen & II & [S]/[E] & St. 2 \\
\hline
6 & einen kurzen Text (3--4 Sätze) über den Klassenraum schreiben & III & [E] & St. 2 \\
\hline
\end{tabular}

% ============================================================================
% 5. METHODISCHE ÜBERLEGUNGEN
% ============================================================================
\section{Methodische Überlegungen}

\subsection{Methodische Großform}

\textbf{Induktive Grammatikvermittlung:}
\begin{enumerate}
    \item Beispiele präsentieren (hay una mesa, hay dos sillas...)
    \item SuS erkennen Muster
    \item Regel gemeinsam formulieren
    \item Anwendung in Übungen
\end{enumerate}

\subsection{Prinzipien}

\begin{itemize}
    \item \textbf{Visualisierung:} Reale Gegenstände + Bildkarten
    \item \textbf{Handlungsorientierung:} Eigene Schultasche beschreiben
    \item \textbf{Kommunikation:} Partnerarbeit mit echtem Informationsaustausch
    \item \textbf{Scaffolding:} Satzanfänge, Wortlisten, Bildunterstützung
\end{itemize}

\subsection{Medien und Materialien}

\begin{tabular}{|l|l|}
\hline
\textbf{Medium} & \textbf{Einsatz} \\
\hline
Tafel/Whiteboard & Tafelbild, Beispielsätze \\
Lehrwerk (S. 36--38) & Vokabel-Input, Übungen \\
Bildkarten (laminiert) & Vokabel-Einführung \\
Reale Gegenstände & Authentizität, Motivation \\
Arbeitsblatt (AB-01-V2) & Festigung, Hausaufgabe \\
Lernpfad (LP-01-V2) & Digitale Übung (optional) \\
\hline
\end{tabular}

% ============================================================================
% 6. VERLAUFSPLANUNG
% ============================================================================
\section{Verlaufsplanung}

\subsection{1. Stunde (45 Min.) -- Vokabeln \& hay entdecken}

\begin{longtable}{|p{0.8cm}|p{1.3cm}|p{4.2cm}|p{3.2cm}|p{1.5cm}|p{1.8cm}|}
\hline
\textbf{Zeit} & \textbf{Phase} & \textbf{Lehrerhandeln} & \textbf{Schülerhandeln} & \textbf{SF} & \textbf{Material} \\
\hline
\endhead

0--3 & Einstieg & Begrüßung: \textit{¡Buenos días!} Zeigt auf Gegenstände: \textit{¿Qué es esto?} & Vermuten (deutsch erlaubt) & Plenum & Raum \\
\hline

3--12 & Input I & Führt 8 Vokabeln ein: zeigt Gegenstand, spricht vor, SuS sprechen nach & Hören, chorisch + einzeln nachsprechen & Plenum & Gegenstände, Tafel \\
\hline

12--18 & Übung I & \textit{¿Qué es?} -- Zeigt auf Gegenstände, SuS benennen & Benennen: \textit{Es una mesa.} & Plenum & Raum \\
\hline

18--25 & Input II & Schreibt Sätze an Tafel: \textit{En la clase hay una mesa. Hay cinco sillas.} Fragt: Was fällt auf? & Beobachten, Muster erkennen, Regel formulieren & Plenum & Tafel \\
\hline

25--35 & Übung II & Kettenübung: LK beginnt \textit{En mi clase hay...}, nächster wiederholt + ergänzt & Kette fortführen, auswendig wiederholen & Plenum & -- \\
\hline

35--42 & Sicherung & Hefteintrag: Vokabeln + hay-Regel (Tafel abschreiben) & Notieren ins Heft & EA & Heft, Tafel \\
\hline

42--45 & Über\-leitung & Ausblick 2. Stunde: \textit{¿Qué hay en tu mochila?} & Hören & Plenum & -- \\
\hline
\end{longtable}

\subsection{2. Stunde (45 Min.) -- Anwendung \& Festigung}

\begin{longtable}{|p{0.8cm}|p{1.3cm}|p{4.2cm}|p{3.2cm}|p{1.5cm}|p{1.8cm}|}
\hline
\textbf{Zeit} & \textbf{Phase} & \textbf{Lehrerhandeln} & \textbf{Schülerhandeln} & \textbf{SF} & \textbf{Material} \\
\hline
\endhead

0--5 & Warm-up & Schnelles Vokabelspiel: Bildkarte hochhalten, SuS rufen Wort & Wort rufen (Wettbewerb) & Plenum & Bildkarten \\
\hline

5--8 & Demo & Öffnet eigene Tasche: \textit{En mi mochila hay un libro, hay dos bolígrafos...} & Zuhören, verstehen & Plenum & Tasche \\
\hline

8--18 & Anwen\-dung I & \textit{Ahora vosotros.} Partnerarbeit: \textit{¿Qué hay en tu mochila?} & Dialog führen, Gegenstände zeigen & PA & Schultaschen \\
\hline

18--22 & Plenum & Einige Paare präsentieren & Präsentieren: \textit{En la mochila de [Name] hay...} & Plenum & -- \\
\hline

22--25 & Erwei\-terung & Einführung: \textit{¿Cuántos/as ... hay?} mit Beispielen & Fragen + Antworten verstehen & Plenum & Tafel \\
\hline

25--37 & Festigung & Verteilt Arbeitsblatt, erklärt Aufgaben, unterstützt individuell & AB bearbeiten (Aufg. 1--4) & EA & AB-01-V2 \\
\hline

37--42 & Sicherung & Lösungen an der Tafel, Nachfragen klären & Vergleichen, korrigieren & Plenum & Tafel \\
\hline

42--45 & Abschluss & Zusammenfassung, HA: Vokabeln lernen, optional LP & Notieren & Plenum & -- \\
\hline
\end{longtable}

\textbf{Legende:} EA = Einzelarbeit, PA = Partnerarbeit, SF = Sozialform, LK = Lehrkraft

% ============================================================================
% 7. TAFELBILD
% ============================================================================
\section{Tafelbild}

\begin{center}
\fbox{\parbox{14cm}{
\textbf{\large Objetos de clase -- hay}

\vspace{0.5em}
\rule{14cm}{0.4pt}
\vspace{0.5em}

\begin{tabular}{p{6cm}|p{6cm}}
\textbf{Vokabeln} & \textbf{hay = es gibt} \\
\hline
la mesa -- der Tisch & \\
la silla -- der Stuhl & En la clase \textbf{hay una} mesa. \\
la pizarra -- die Tafel & \textbf{Hay cinco} sillas. \\
la ventana -- das Fenster & \textbf{Hay muchos} libros. \\
la puerta -- die Tür & \\
el libro -- das Buch & \rule{6cm}{0.4pt} \\
el cuaderno -- das Heft & \\
el bolígrafo -- der Kuli & \textbf{¿Qué hay en la clase?} \\
& $\rightarrow$ Hay una pizarra. \\
\textbf{Achtung:} el mapa (m.) & \\
\end{tabular}

\vspace{0.5em}
\rule{14cm}{0.4pt}
\vspace{0.3em}

\textit{hay + un/una} \quad | \quad \textit{hay + Zahl} \quad | \quad \textit{hay + muchos/as}

\vspace{0.3em}
\textbf{NICHT:} *hay el libro \quad *hay la mesa
}}
\end{center}

% ============================================================================
% 8. ERWARTETE SCHWIERIGKEITEN
% ============================================================================
\section{Erwartete Schwierigkeiten und Reaktionen}

\begin{tabular}{|p{4.5cm}|p{4.5cm}|p{4.5cm}|}
\hline
\textbf{Schwierigkeit} & \textbf{Ursache} & \textbf{Reaktion} \\
\hline
*Hay el libro & Transfer aus Deutsch (``Es gibt das Buch'') & Kontrastübung: hay = neu/unbekannt \\
\hline
Aussprache \textit{ll} (silla) & Keine Entsprechung im Deutschen & Vorsprechen, chorisch üben: [ʝ] \\
\hline
el mapa (maskulin) & Erwartung: -a = feminin & Merkspruch: ``Der Mapa ist ein Mann!'' \\
\hline
Verwechslung hay/es & Englisch-Einfluss (there is) & Klare Abgrenzung: hay $\neq$ es \\
\hline
Zeitproblem & Heterogene Arbeitsgeschwindigkeit & Puffer: Kettenübung flexibel, AB kürzen \\
\hline
Lustlosigkeit bei Vokabeln & Zu abstrakt & Echte Gegenstände, Schultaschenspiel \\
\hline
\end{tabular}

% ============================================================================
% 9. DIFFERENZIERUNG IM DETAIL
% ============================================================================
\section{Differenzierung}

\subsection{Während der Stunde}

\begin{tabular}{|l|p{4cm}|p{4cm}|p{4cm}|}
\hline
\textbf{Phase} & \textbf{[B] Basis} & \textbf{[S] Standard} & \textbf{[E] Erweiterung} \\
\hline
Vokabel-Input & Bildkarten als Stütze & Ohne Hilfe benennen & Zusatzvokabeln selbst nachschlagen \\
\hline
hay-Übung & Lückensätze ausfüllen & Sätze frei bilden & Mini-Text: 4 Sätze \\
\hline
Partnerarbeit & 3 Gegenstände beschreiben & 6 Gegenstände & + Nachfragen: ¿Cuántos...? \\
\hline
Arbeitsblatt & Aufgaben 1--2 & Aufgaben 1--4 & Aufgaben 1--5 + Bonus \\
\hline
\end{tabular}

\subsection{Scaffolding-Materialien}

\begin{itemize}
    \item \textbf{[B]:} Wortkarten mit Bild + Artikel (zum Abschauen)
    \item \textbf{[S]:} Satzanfänge: \textit{En mi clase hay... / En mi mochila hay...}
    \item \textbf{[E]:} Impulskarte: \textit{Describe tu clase ideal.}
\end{itemize}

% ============================================================================
% 10. HAUSAUFGABE
% ============================================================================
\section{Hausaufgabe}

\begin{itemize}
    \item \textbf{Pflicht:} Vokabeln lernen (8 Kernwörter)
    \item \textbf{Pflicht:} Lehrwerk S. 37, Aufgabe 2 (schriftlich)
    \item \textbf{Optional:} Lernpfad LP-01-V2 bearbeiten (digital)
    \item \textbf{[E]:} 5 Sätze über das eigene Zimmer: \textit{En mi habitación hay...}
\end{itemize}

% ============================================================================
% 11. REFLEXIONSFRAGEN
% ============================================================================
\section{Reflexionsfragen (für Lehrkraft)}

\begin{enumerate}
    \item Konnten alle SuS mindestens 6 Vokabeln aktiv benennen?
    \item War die induktive Regelableitung (hay) für alle verständlich?
    \item Wie war das Verhältnis LK-Sprechzeit zu SuS-Sprechzeit?
    \item War die Kettenübung motivierend oder zu lang?
    \item Reichte die Zeit für das Arbeitsblatt?
    \item Welche Fehler traten gehäuft auf? (für nächste Stunde)
\end{enumerate}

% ============================================================================
% 12. ANLAGEN
% ============================================================================
\section{Anlagen}

\begin{enumerate}
    \item \textbf{AB-01-mi-instituto-V2.pdf} -- Arbeitsblatt (Klassensatz)
    \item \textbf{ML-01-mi-instituto-V2.pdf} -- Musterlösung
    \item \textbf{LP-01-mi-instituto-V2.html} -- Digitaler Lernpfad
    \item Bildkarten: Klassenzimmer (laminiert, 12 Stück)
    \item Wortkarten: Spanisch-Deutsch (laminiert)
\end{enumerate}

% ============================================================================
% 13. DIDAKTIK-SCORE
% ============================================================================
\newpage
\section{Didaktik-Score}

\begin{scorebox}
\textbf{Bewertung der didaktischen Qualität nach 10 Dimensionen}\\
\textbf{Ziel: $\geq$ 9.0 (Freigabe)}

\vspace{1em}
\begin{tabular}{|c|l|c|c|p{4.5cm}|}
\hline
\textbf{\#} & \textbf{Dimension} & \textbf{Gew.} & \textbf{Score} & \textbf{Begründung} \\
\hline
1 & Differenzierung & 15\% & 9/10 & [B]/[S]/[E] qualitativ verschieden, Scaffolding \\
\hline
2 & Taxonomie (AFB) & 10\% & 9/10 & AFB I: 30\%, II: 50\%, III: 20\% \\
\hline
3 & Lernziel-Alignment & 10\% & 10/10 & Alle FZ aus Rahmenplan MV ableitbar \\
\hline
4 & Aktivierung & 10\% & 9/10 & Hoher SuS-Sprechanteil, Kettenübung, PA \\
\hline
5 & Scaffolding & 10\% & 9/10 & Bildkarten, Satzanfänge, gestufte Hilfen \\
\hline
6 & Feedback-Qualität & 10\% & 8/10 & Chorisch + individuell, aber wenig elaboriert \\
\hline
7 & Sprachliche Klarheit & 10\% & 10/10 & Kurze Sätze, altersgerecht, klare Instruktionen \\
\hline
8 & Motivation & 10\% & 9/10 & Eigene Schultasche = Alltagsbezug \\
\hline
9 & Inklusion (UDL) & 10\% & 9/10 & Visuell + auditiv + haptisch \\
\hline
10 & Praktikabilität & 5\% & 9/10 & 90 Min. realistisch, Material vorhanden \\
\hline
\multicolumn{3}{|r|}{\textbf{GESAMT:}} & \textbf{9.1/10} & \textcolor{successgreen}{\textbf{FREIGABE}} \\
\hline
\end{tabular}

\vspace{1em}
\textbf{Status:}
\begin{itemize}
    \item[$\boxtimes$] \textcolor{successgreen}{$\geq$ 9.0: \textbf{FREIGABE} (Premium-Qualität)}
    \item[$\Box$] 8.0--8.9: Iteration empfohlen
    \item[$\Box$] $<$ 8.0: Überarbeitung erforderlich
\end{itemize}

\vspace{0.5em}
\textbf{Verbesserungspotential:}
\begin{itemize}
    \item Dimension 6 (Feedback): Mehr elaboriertes Feedback einplanen (z.B. Peer-Korrektur)
\end{itemize}
\end{scorebox}

\end{document}
